\documentclass{report}
\usepackage{amsmath,amssymb}
\setlength{\parindent}{0mm}
\setlength{\parskip}{1em}
\begin{document}
\begin{center}
\rule{6in}{1pt} \
{\large Stephen Nash \\
{\bf Using Inexact Gradients in a Multilevel Optimization Algorithm}}

Volgenau School of Engineering \\ George Mason University \\ Nguyen Engineering Building \\ Room 2500 \\ Mailstop 5C8 \\ Fairfax \\ VA 22030
\\
{\tt snash@gmu.edu}\\
Michael Lewis\end{center}

Many optimization algorithms require gradients of the model functions,
but computing accurate gradients can be computationally expensive. We
study the implications of using inexact gradients in the context of the
multilevel optimization algorithm MG/OPT. MG/OPT recursively uses
(typically cheaper) coarse models to obtain search directions for
finer-level models. However, MG/OPT requires the gradient on the fine
level to define the recursion. Our primary focus here is the impact of
the gradient errors on the multilevel recursion. We analyze, partly
through model problems, how MG/OPT is affected under various assumptions
about the source of the error in the gradients, and demonstrate that in
many cases the effect of the errors is benign. Computational experiments
are presented.


\end{document}
